% ============================================================
% OP_anomaly_proposal_v1.tex
% Proposal: anomaly-cancellation package (OP-GUT6 family)
% For GPT review. NO edits to the preprint are made by this file.
% ============================================================
\section*{Proposal: anomaly-cancellation package (OP-GUT6a sharpening)}

\subsection*{0. Scope and relation to existing material}
Existing infrastructure (no duplication intended):
\S\ref{sec:anomaly_architecture} (Part~I) already contains the four
local conditions (eqs.~anom\_3Y, anom\_2Y, anom\_YYY, anom\_Ygrav),
the inherited assignments (eq.~hypercharges), the hard disclaimer B3,
the prose-level statement B4 ($\mathcal L_5$ anomaly-transparent),
the Witten counting B5, Layers C1--C4, and the status table.
\S\ref{sec:anomaly_constraints} (Part~IV) restates the constraints
under P7 and asserts ``these constraints fix hypercharges up to
normalisation'' for the one-generation content.
The charge-quantisation block (\S\ref{sec:common_origin}) frames
anomaly freedom as a mandatory self-consistency requirement and
records the lattice constraint $2T+Y\in2\mathbb Z$.
Registry: OP-GUT6a (local triangles), OP-GUT6b (Witten),
OP-GUT6c (emergent-graviton interpretation), all Open.

\emph{Gaps this package addresses:}
(i) the conditions are stated but the arithmetic is never displayed;
(ii) the Part-IV uniqueness sentence is asserted without proof or
caveats, although the actual solution structure has a degenerate
branch and, with an uncharged-singlet freedom, an exact flat
direction;
(iii) B4 is a prose claim that can be upgraded to a short lemma.

\subsection*{1. Decomposition (mirroring OP3 / OP-Q11 style)}
\begin{itemize}
  \item \textbf{AN1} --- explicit arithmetic verification that the
    inherited assignments satisfy all four conditions
    (candidate Level~A arithmetic of \emph{inherited} content;
    verification, not derivation).
  \item \textbf{AN2} --- lemma: $\mathcal L_5$ is anomaly-transparent
    (candidate Level~A at the one-loop/anomaly-polynomial level;
    ECT application inherits B1 conditionality).
  \item \textbf{AN3} --- classification of the solution set of the
    four conditions for the one-generation content (candidate
    Level~A arithmetic, conditional on inherited representation
    content); corollaries: precise scoping of the Part-IV
    ``fixed up to normalisation'' sentence; the $\nu^c$ flat
    direction; what compactness does and does not add.
  \item \textbf{AN4} --- unchanged open remainder: OP-GUT6c,
    't~Hooft matching (C4), dynamical derivation of matter content
    (OP-GUT1/2/3b/5). No new claims.
\end{itemize}
OP-GUT6a/6b registry rows stay \emph{Open as derivations}; comments
gain ``explicit consistency verification recorded for the inherited
content''.

\subsection*{2. AN1: explicit verification (proposed text core)}
With eq.~hypercharges
($Y_Q=\tfrac16$, $Y_{u^c}=-\tfrac23$, $Y_{d^c}=\tfrac13$,
$Y_L=-\tfrac12$, $Y_{e^c}=1$, optionally $Y_{\nu^c}=0$):
\begin{align*}
[SU(3)]^2U(1)_Y&:\ 2\cdot\tfrac16-\tfrac23+\tfrac13=0,\\
[SU(2)]^2U(1)_Y&:\ 3\cdot\tfrac16-\tfrac12=0,\\
U(1)_Y\text{-grav}^2&:\ 6\cdot\tfrac16-3\cdot\tfrac23+3\cdot\tfrac13
  -2\cdot\tfrac12+1\;(+0)=1-2+1-1+1=0,\\
[U(1)_Y]^3&:\ 6(\tfrac16)^3+3(-\tfrac23)^3+3(\tfrac13)^3
  +2(-\tfrac12)^3+1^3\;(+0)\\
 &\quad=\tfrac1{36}-\tfrac{32}{36}+\tfrac{4}{36}-\tfrac{9}{36}
  +\tfrac{36}{36}=0.
\end{align*}
Witten: $4N_{\rm gen}$ doublets, even for all integer
$N_{\rm gen}$ (restates B5 numerically).
Status: candidate Level~A \emph{arithmetic of inherited input};
explicitly not a derivation; does not change the B3 scope labels
(which conditions are \emph{meaningful} in current ECT is a separate
classification from whether the inherited template values satisfy
them).

\subsection*{3. AN2: lemma ($\mathcal L_5$ anomaly transparency)}
\emph{Lemma (proposed).}
In the reconstructible ordered-branch chiral gauge EFT, the
deformation $\mathcal L_5=\mu_5\bar\Psi\gamma^\mu n_\mu\Psi$ leaves
every gauge-anomaly coefficient of the anomaly polynomial, and the
Witten $\mathbb Z_2$ class, unchanged.
\emph{Proof sketch.}
(i) The anomaly coefficients are functions of the chiral
representation/charge data alone; $\mathcal L_5$ is the coupling of
the background vector $B_\mu=\mu_5 n_\mu$ to the gauge-singlet
\emph{vector} current $\bar\Psi\gamma^\mu\Psi$ and changes no
representation datum.
(ii) Fujikawa form: the chiral Jacobian is computed from the
gauge-covariant Dirac operator; a background \emph{vector} shift
acts identically on both chiralities, so the
$\mathrm{tr}\,\gamma_5\{\cdots\}$ coefficients of the anomaly
polynomial are unchanged at the polynomial level.
(iii) The doublet count is manifestly unchanged (Witten).
\emph{Scope.} One-loop/anomaly-polynomial level; flavour-diagonal
$\mathcal L_5$ as defined; possible 't~Hooft-type bookkeeping
involving the background direction is outside the lemma; ECT
application inherits the B1 OS-closure conditionality.
Status: candidate Level~A (mathematical, at the stated level);
upgrades the existing B4 prose without changing its message.

\subsection*{4. AN3: solution structure of the hypercharge system}
\emph{Proposition (proposed; candidate Level~A arithmetic,
conditional on the inherited one-generation chiral content).}
Impose eqs.~anom\_3Y, anom\_2Y, anom\_Ygrav, anom\_YYY on
$\{Y_Q,Y_{u^c},Y_{d^c},Y_L,Y_{e^c}\}$ (no $\nu^c$).
The linear conditions give $Y_L=-3Y_Q$, $Y_{e^c}=6Y_Q$,
$Y_{u^c}+Y_{d^c}=-2Y_Q$; writing $Y_{u^c}=-Y_Q+\delta$,
$Y_{d^c}=-Y_Q-\delta$, the cubic condition reduces to
\[
18\,Y_Q\,(9Y_Q^2-\delta^2)=0 .
\]
Hence exactly two branches:
(a) $\delta=\mp3Y_Q$: the Standard-Model assignment up to overall
normalisation and the $u^c\!\leftrightarrow\! d^c$ relabelling;
(b) the degenerate branch $Y_Q=0$
($Y_L=Y_{e^c}=0$, $Y_{u^c}=-Y_{d^c}=\delta$ free), under which the
lepton sector is electrically neutral; it is excluded by the
inherited requirement of a charged lepton ($Q_e\neq0$), an input
recorded explicitly as such.
\emph{Corollary 1 (Part-IV scoping).}
The sentence ``these constraints fix hypercharges up to
normalisation'' in \S\ref{sec:anomaly_constraints} is correct for
the listed content $\{Q_L,u_R,d_R,L_L,e_R\}$ \emph{modulo} the
relabelling and the excluded degenerate branch; propose adding this
scoping in one sentence with a pointer to the new material.
\emph{Corollary 2 ($\nu^c$ flat direction).}
If a singlet $\nu^c$ with \emph{free} hypercharge is admitted, the
shifted assignment $Y\to Y+t\,(B{-}L)$ is anomaly-free for all $t$
(all four conditions verified explicitly; standard result), so the
anomaly system alone no longer fixes hypercharge.
The existing Part-IV sentence ``a singlet can be added without
disturbing anomaly balance'' implicitly sets $Y_{\nu^c}=0$; propose
making that explicit.
\emph{Corollary 3 (what compactness adds and does not add).}
The ECT compact-$U(1)$ lattice (Theorem~charge\_lattice;
$2T+Y\in2\mathbb Z$) enforces rationality/lattice structure but does
\emph{not} remove the $B{-}L$ flat direction (which is
rational-valued); the observed assignment therefore remains
OP-GUT3b-level input.
This corrects a tempting overclaim in advance.

\subsection*{5. Implementation map (sites; edits only after approval)}
\begin{enumerate}
  \item \S anomaly\_architecture, Layer~B: new paragraph
    ``B2$'$. Explicit verification for the inherited assignments''
    (AN1 display) after B2.
  \item Layer~B, B4: insert the AN2 lemma + sketch (keep existing
    prose as surrounding context).
  \item Layer~C: new paragraph ``C1$'$. Solution structure of the
    hypercharge system'' (AN3 + corollaries), adjacent to C1.
  \item Status table tab:anomaly\_architecture: add rows for
    AN1 / AN2 / AN3 with candidate-A wording; adjust the
    $\mathcal L_5$ row comment.
  \item \S anomaly\_constraints (Part~IV): one-sentence scoping of
    the uniqueness claim + explicit $Y_{\nu^c}=0$ remark + pointer.
  \item OP-GUT6a/6b rows: comment update (statuses remain Open).
  \item Consistency sweep: grand tables, Part-I derivation graph
    (check for an anomaly node at implementation), symbols/abbrev
    (ABJ presence check), companion (defer, Q11 pattern).
\end{enumerate}

\subsection*{6. Status discipline}
Forbidden: ``ECT derives anomaly cancellation'';
``anomalies (+ compactness) fix the observed hypercharges'';
``ECT predicts the SM matter content''.
Safe: ``verified for the inherited assignments'';
``classification of the solution set under stated assumptions'';
``consistency test passed by the inherited content''.
OP-GUT6a/6b/6c remain Open; OP-GUT3b untouched.

\subsection*{7. References to add (data to be page-verified before
insertion --- Kikukawa lesson)}
Candidates: Geng--Marshak (1989); Minahan--Ramond--Warner (1990);
optionally a review (Foot--Lew--Volkas).
Witten1982 and AharonyEtAl2013 already in the bibliography.

\subsection*{8. Questions for GPT}
(Q1) Approve the AN1/AN2/AN3 decomposition and candidate-A wording?
(Q2) AN3 placement: Layer~C paragraph vs.\ separate subsubsection?
(Q3) Is the $Q_e\neq0$ exclusion of the degenerate branch the
cleanest inherited input, or prefer ``all SM fermions chiral under
$U(1)_{\rm em}$''-type wording?
(Q4) Lemma AN2: is the Fujikawa sketch sufficient at this status
level, or restrict the statement to the representation-data argument
only?
(Q5) Which references to include; any objection to citing the
classification literature for AN3?
